% Section III -- Design (SKELETON; text written in stage C after GATE-RT).
\section{Design}\label{sec:design}

\subsection{Architecture}
\sremustub{Processes and planes (F1): TX and RX applications, external controller
processes, engine, native framed ZMQ and raw REQ/REP data planes versus the JSON control
plane, recorder (SigMF, events.jsonl, manifest). DESIGN sections 1.1, 8.}
\begin{figure*}[t]
\centering
\sremufigure{F1}{\textwidth}{Architecture of \sremu{}: applications, controllers, engine,
data planes, control plane and recorder.}{fig:F1}
\end{figure*}

\subsection{Time model and command semantics}
\sremustub{Exact rational time base and rounding rules (R1--R5); page and segment grid;
switch time defined at the RIS and applied per receiver at to\_ticks of the switch time
plus the RIS-to-receiver delay; late commands; same-time ties; transitions; lockstep
barriers and paced mode (F2). DESIGN sections 4.1--4.10.}
\begin{figure*}[t]
\centering
\sremufigure{F2}{\textwidth}{(a) Page and segment processing. (b) Life cycle of a RIS
command from wall-clock receipt to the applied sample on each receiver, with one late
command and one barrier.}{fig:F2}
\end{figure*}

\subsection{Interfaces}
\sremustub{UHD-shaped Python API, compatibility wrappers, native ZMQ frame header, control
socket vocabulary, security defaults; Table III maps sremu calls to UHD, srsRAN ZMQ and
GNU Radio counterparts (cite UHDManual, srsRANZMQ, SigMF). DESIGN section 8.}

\subsection{State library and ETSI GS RIS 008}
\sremustub{Response-table data model, provenance, load-time validators, importers and
built-in tables; Table VI maps GR RIS 002/003 and GS RIS 008 parameters to sremu elements,
worded as alignment. DESIGN sections 5.5, 8.9.}
